% Standalone Overleaf project. Compiler: pdfLaTeX. Main document: main.tex.
% S4.1 is measured; the S4.2 section is a prospective proposal, not executed results.
\documentclass[11pt,a4paper]{article}
\usepackage[margin=2.15cm]{geometry}
\usepackage[T1]{fontenc}
\usepackage[utf8]{inputenc}
\usepackage{lmodern,microtype,amsmath,amssymb,booktabs,array,longtable,graphicx}
\usepackage{enumitem,float,placeins,xcolor}
\usepackage[hidelinks]{hyperref}
\setlength{\parindent}{0pt}
\setlength{\parskip}{5pt}
\setlength{\emergencystretch}{3em}
\setlist{nosep,leftmargin=1.5em}
\newcommand{\file}[1]{\nolinkurl{#1}}
\newcommand{\E}{\mathbb{E}}
\title{Experiment Stage 4: Causal Communication and RI Participation\\
\large S4.1 results and an evidence-based S4.2 measurement plan}
\author{Single-hop in-context mother-of reasoning\\OLMo-2-1124-7B}
\date{24 September 2026}
\begin{document}
\maketitle

\begin{abstract}
The completed S4.1 run identifies a reproducible set of answer-affecting routes, but does not yet establish a jointly sufficient circuit. L17H1 transmits substantial fact-site influence through the V channels of L18H18 and L18H19; L15H25 reaches L16H1 and L16H21 through V at child-last. Colon-source effects from the latter pair reach the former pair through Q, suggesting a convergence motif that requires joint tests. Additional routes have consistent negative effects and must not be removed from the explanation merely because of their sign. Of 96 selected routes extended to 89 discovery families, 95 satisfy the operational retention rule under noising and 93 under restoration; on the 69 families outside the initial mapping panel, the corresponding counts are 94 and 92. Strong effects are sensitive to fact order, and several other-fact controls are non-negligible. Coverage also identifies L13H18 as a causally relevant source whose tested direct routes remain unresolved. All 70,408 stored records, code/input identities, adaptive decisions and exported analyses were verified locally. This report separates measured routes from proposed groups, and specifies a deduplicated S4.2 plan without claiming results from that future stage.
\end{abstract}

\setcounter{tocdepth}{1}
\tableofcontents
\newpage
\section{Scope, provenance and verification}
\label{sec:verification}
This report analyzes \file{stage4_all_v2}, Slurm job 923239, using the executed v2 package and actual adaptive schedules. The pinned model revision is
\file{7df9a82518afdecae4e8c026b27adccc8c1f0032}. The recorded runtime is PyTorch 2.14.0+cu126 and Transformers 4.51.3 on node s-004. Six GPUs were arranged as three independent two-GPU workers. No new model run was performed for this analysis.

The source archive SHA-256 is:
\par{\footnotesize\file{30d3ed8dbb95808c426b7d8b3faa4f437b9962b9872a17a7dd8f63689848ac5f}}\par
The downloaded package's declared hashes were checked against available package files and, for shared dependencies, exact matching local sources. All worker code hashes agree with the controller manifest. Every checksummed chunk was verified, and every expected pair/configuration/direction identifier was present exactly once in the appropriate phase. There were no unmarked JSON chunks. The original outputs were preserved.

\begin{table}[htbp]\centering\small
\begin{tabular}{lrrr}
\toprule
Phase & Configs & Records & Reused \\
\midrule
coverage & 4 & 712 & 80 \\
localize\_roles & 6 & 240 & 0 \\
localize\_slots & 8 & 320 & 0 \\
seed & 769 & 30,760 & 0 \\
refine & 105 & 4,200 & 0 \\
extend\_seed & 72 & 25,632 & 0 \\
extend\_refine & 24 & 8,544 & 0 \\
\bottomrule
\end{tabular}

\caption{Actual completed measurement schedule. Record counts include both directions in the extension phases. They exclude gate and capture forwards.}
\end{table}

There are 70,408 records: 70,328 newly measured endpoints and 80 reused Stage-2 Scope-P effects, explicitly labelled as reused. The 584-configuration seed preview was expanded by the measured coverage decisions to 769 configurations. All 21 measurement workers and the initial gate worker passed six-pair full-model gates and the mandatory tiny-model tests. Across saved gates, no-op, temporal-null, cross-position-null, QKV/output equivalence and saved patch/logit errors were zero; the maximum recorded SDPA reconstruction error was $6.10\times10^{-5}$. The 2,856 recorded per-configuration self-endpoint checks also had zero drift. These are checks of the implemented interventions, not evidence that every scientific hypothesis is correct.

All four exported CSV tables in each of the seven measurement phases were regenerated and matched within $10^{-12}$ absolute/relative tolerance. The two interaction tables were reproduced, and all six adaptive scheduling decisions were reconstructed with identical scientific content. An independent arithmetic check recovered each new endpoint effect from its two answer logits. These checks establish internal integrity and reproducibility; they do not independently validate every modeling assumption.

\section{What was measured}
\label{sec:measurement}
There are 89 discovery families, each in two orders (178 paired clean/corrupted inputs). Initial mapping uses the fixed 20-family/40-pair subset. The remaining 69 families are an expansion check, \emph{not} a pristine confirmatory population: earlier research stages already used them. The 87 held-out families were not opened.

With a fixed clean-answer sign, the metric is
\[
 M(x)=z(y_c\mid x,p)-z(y_r\mid x,p),\qquad
 I=M(x_c)-M(x_c;T_{r\to c}),\quad
 J=M(x_r;T_{c\to r})-M(x_r).
\]
The shared answer prefix $p$ is teacher-forced and the first divergent answer token is scored. Positive $I$ means noising harms the clean-answer margin; positive $J$ means restoration improves it. Negative effects are preserved. Orders are averaged within family before equally weighting families. The measured full clean-minus-corrupted gap is 11.17046 logits. Intact candidate-margin accuracy is 178/178 in both conditions on this population; this is not a full-name generation metric.

\textbf{A route is an answer intervention, not just an activation difference.}
The source head's pre-output-projection $z=AV$ slice is replaced with its donor value. Its shared attention-output normalization stays live. Other intervening residual branch increments are frozen \emph{after} normalization. The resulting receiver channel is captured and injected in a fresh recipient run, with subsequent computation live. For a fact-site K/V route, injection occurs at the fact positions, but only the receiver head's original colon output row is released. For Q, injection is at the colon. MLP tests inject the changed local MLP input; bypass tests keep downstream branches frozen and the final normalization live.

Thus a measured route is conditional on its declared recipient background and live/frozen boundaries. It does not prove a semantic code in an isolated head, establish an additive fraction of total importance, or prove that independently measured edges compose into an end-to-end chain. In particular, shared normalization prevents attributing a route to a linearly isolated normalized head contribution.

\textbf{Fixed retention rule.} With $\tau=0.10$ logits, a coherent candidate has $|\bar I|\ge\tau$ and at least 70\% family sign agreement. A heterogeneous candidate has $\E_f|I_f|\ge\tau$ and at least 20\% of families with $|I_f|\ge\tau$. The thresholds prioritize measurements; they are not significance tests. We retain signed and absolute effects, direction, per-family variation and 20,000-draw family-bootstrap intervals (seed 20260923). All discovery intervals are descriptive after selection. No multiple-comparison-confirmed membership claim is made here.

\section{Coverage changes the candidate map}
\label{sec:coverage}
\begin{table}[htbp]\centering\small
\begin{tabular}{llrrr}
\toprule
Head & Scope & Mean $I$ & Mean $|I_f|$ & Sign fraction \\
\midrule
L13H18 & all & 0.314 & 0.417 & 0.80 \\
L13H18 & colon & 0.022 & 0.027 & 0.81 \\
L24H19 & all & 0.344 & 0.346 & 0.97 \\
L24H19 & colon & 0.349 & 0.350 & 0.98 \\
\bottomrule
\end{tabular}

\caption{Outside-inventory candidates on all 89 discovery families. Scope P covers original prompt positions; Scope F is the original colon.}
\end{table}

\textbf{L13H18 is a supported unresolved fact-site candidate.}
Its P effect is $+0.314$, compared with F $+0.022$. On the common panel, the joint all-mother mask has effect $+0.230$ and the all-template mask $+0.149$. Slot localization retains fact1-mother, fact3-mother and fact3-template (zero-based slot names). These are jointly measured masks, not sums of token effects. Slot0 masks have zero donor effects in these measurements; this is not a general statement of positional irrelevance.

The resulting 138 direct configurations from the three retained source masks (132 head-channel and six MLP tests) yield no retained route. The largest mean absolute family effect is only $0.035$, for fact3-template to L18H18 K. This supports a bounded local mediation follow-up in S4.3. It does not exclude L13H18, prove that it is redundant, or justify adding it to a specific G1 consumer group without evidence.

\textbf{L24H19 is principally an answer-position candidate.}
Its P and F effects are $+0.344$ and $+0.349$. Its extended colon residual bypass has $I=+0.165$, while no tested head/MLP outgoing route meets the common-panel rule. The bypass is a real isolated residual effect under the frozen background, but cannot be subtracted from the total F effect to estimate an unexplained causal remainder. Both candidates are non-RI coverage additions; their distinct outcomes illustrate why coverage should not be inferred from the original shortlist.

\section{The measured communication map}
\label{sec:map}
Of 769 initial configurations, 76 meet the common-panel rule: 62 head-channel routes, seven MLP routes and seven bypasses. There are 115 retained non-control configurations after combining seed and refinement; 96 were extended and 19 remain common-panel only because of the fixed budget. Multiple configurations can refer to the same head pair with different sites/channels and must not be counted as distinct circuit members.

\begin{figure}[htbp]\centering
\includegraphics[width=\linewidth]{figures/route_schematic.png}
\caption{A selected communication motif, with 89-family mean effects. Each arrow summarizes separately measured routes; grouped boxes do not denote a measured joint intervention. The full ledger includes additional and opposing routes. In particular, the displayed successive arrows have not yet been tested as a complete mediated chain.}
\label{fig:map}
\end{figure}

\subsection{L17H1: strong V routes with both effect signs}
The clearest fact-to-answer effects run from L17H1 into L18H18 and L18H19 V. Combining query-mother, \texttt{is} and period as an explicitly measured union gives $I=+3.400$ and $+2.074$, with $J=+3.335$ and $+2.054$. Both noising effects have the positive sign in all 89 families. The additional-69 means are $+3.391$ and $+2.057$, close to the initial-panel values.

The same union has negative routes into L18H24 ($-1.004$), L20H1 ($-0.582$), and L19H16 ($-0.247$), and a positive route into L22H5 ($+0.898$). The two leading positive means are 30.4\% and 18.6\% of the full-model gap on the same population; these are scale-normalized effects, not explained fractions, and must not be added as mediation shares. The route ledger provides the corresponding ratio of means for every extended route. Consequently, a common source supports routes with opposite effects on the answer contrast. Group tests must retain that sign structure. Summing the positive routes would overstate an unmeasured joint contribution; removing negative routes merely to increase a performance score would also distort the mechanism.

K-only routes from this source did not meet the initial retention rule. This supports V as the leading \emph{tested direct channel}, not a claim that keys never matter. Joint KV tests and their interactions below show why a small standalone K effect does not establish separability.

\begin{figure}[htbp]\centering
\includegraphics[width=\linewidth]{figures/fact_routes.png}
\caption{Selected receiver columns of the complete common-panel V measurements. Cells contain signed mean effects, including below-threshold effects. Blank cells are temporally ineligible, not measured zeros. Different source sites are separate interventions.}
\end{figure}

\subsection{L15H25: a narrower route into the layer-16 pair}
At child-last, L15H25 reaches L16H1 and L16H21 through V with extended effects $+0.182$ and $+0.210$; reverse restoration gives $+0.187$ and $+0.223$. Joint KV is very similar ($+0.185$ and $+0.212$). No local MLP route meets the initial rule. These results give a concrete consumer hypothesis despite the inconclusive Stage-3 readout; they do not determine whether the transmitted representation is an address, a name identity, or another feature.

The layer-16 heads in turn have colon Q routes into L18H18 and L18H19. Together with L17H1's fact-site V routes, this suggests a convergence of fact-site influence and colon-state influence on the same pair. It is premature to call Q the query-selection code: this run used fact-swap donors, not a query-only intervention.

\subsection{Colon routes: downstream continuation and parallel exits}
L18H18 and L18H19 have Q routes into L27H6 ($+0.662$ and $+0.731$), KV/V routes into L21H23, and additional routes toward L25H17 and several negative contributors. They also have measured MLP19 routes ($+0.410$ and $+0.532$) and residual bypasses ($+0.547$ and $+0.754$). Hence their answer effects are not supported only through a single type of downstream exit. These routes overlap causally in the live network and are not an additive partition.

Negative sources also have supported downstream routes: L20H1, for example, reaches L21H23 KV and L27H6 Q with negative mean effects. A negative route from a positive source to a negative head, or vice versa, must be described by the measured source/site/channel/background; signs are not interchangeable with a global ``good head/bad head'' label.

\subsection{Secondary sources and gaps}
L13H10 at the whole query sentence reaches L16H1 V negatively ($-0.175$), L17H3 V positively ($+0.150$), L18H18 K negatively ($-0.107$), L18H19 V negatively ($-0.116$), and MLP14 positively ($+0.184$). This mixed pattern argues against assigning one universal sign to all of its outgoing routes. L14H26 has one retained head recipient, L16H1 V ($+0.136$); it does not yet justify a multi-member outgoing group.

L17H3's retained routes originate at the colon and reach L18H18/L18H19 Q negatively ($-0.158/-0.106$). Its early masks yield no retained route. This strengthens the colon/mixed interpretation and does not establish a strong early writer.

Some previously important heads remain weakly explained by this map, notably L21H18 and L21H6. L26H23's prior causal relevance and L26H31's backup hypothesis are also not resolved by this direct-route panel. These are explicit coverage boundaries, not negative membership results.

\begin{table}[htbp]\centering\small
\begin{tabular}{llrrr}
\toprule
Source/site & Receiver/channel & $I_{89}$ & $J_{89}$ & $I_{69}$ \\
\midrule
L17H1/union & L18H18/V & +3.400 & +3.335 & +3.391 \\
L17H1/union & L18H19/V & +2.074 & +2.054 & +2.057 \\
L17H1/union & L18H24/V & -1.004 & -1.070 & -0.988 \\
L17H1/union & L22H5/V & +0.898 & +0.865 & +0.882 \\
L17H1/union & L20H1/V & -0.582 & -0.576 & -0.567 \\
L17H1/union & L19H16/V & -0.247 & -0.225 & -0.243 \\
L15H25/child-last & L16H1/V & +0.182 & +0.187 & +0.182 \\
L15H25/child-last & L16H21/V & +0.210 & +0.223 & +0.215 \\
L16H1/colon & L18H18/Q & +0.372 & +0.370 & +0.367 \\
L16H21/colon & L18H18/Q & +0.349 & +0.407 & +0.348 \\
L18H18/colon & L27H6/Q & +0.662 & +0.663 & +0.674 \\
L18H19/colon & L27H6/Q & +0.731 & +0.742 & +0.736 \\
L13H10/sentence & L16H1/V & -0.175 & -0.175 & -0.172 \\
L13H10/sentence & L17H3/V & +0.150 & +0.139 & +0.151 \\
L14H26/sentence & L16H1/V & +0.136 & +0.140 & +0.139 \\
L17H3/colon & L18H18/Q & -0.158 & -0.172 & -0.155 \\
L18H19/colon & MLP 19 & +0.532 & +0.554 & +0.514 \\
L18H19/colon & bypass & +0.754 & +0.770 & +0.748 \\
L18H18/colon & bypass & +0.547 & +0.562 & +0.546 \\
L24H19/colon & bypass & +0.165 & +0.165 & +0.160 \\
\bottomrule
\end{tabular}

\caption{Selected extended routes in logits. $I_{89}$ and $J_{89}$ average 89 family means; $I_{69}$ uses only the additional families. ``Union'' means query-mother, \texttt{is} and period; ``sentence'' is the query fact. Numerical details, intervals and all 96 configurations are provided in \file{data/route_ledger_96.csv} and \file{data/extension_population_summary.csv}.}
\label{tab:routes}
\end{table}
\FloatBarrier

\section{Expansion, controls and interaction evidence}
\label{sec:controls}
\subsection{Expansion is largely consistent, with informative boundary cases}
The pipeline selected 72 seed routes and 24 refinement routes for both-direction measurement on all 178 pairs. On 89 families, noising retains 93 coherent and two heterogeneous routes; restoration retains 91 coherent and two heterogeneous routes. Mean signs agree between noising and restoration for all 96 routes, although signs alone do not imply equal magnitudes or strong effects. The 69-family panel retains 94 noising and 92 restoration routes.

L17H24 to L21H23 V falls below the rule in both directions on the expanded panel ($I=+0.092$, $J=+0.092$). Its joint KV counterpart was retained on the common panel but not extended; failure of V alone is not a measured KV failure. L18H18 to L22H5 Q is borderline ($I=+0.099$, mean absolute effect $0.100$) and drops below the rule in the additional families and in restoration. L18H19 to MLP18 also fails the restoration rule. The L17H1 mother-to-MLP17 route remains heterogeneous on all families, but its reverse effect does not retain on the additional 69 alone. These cases must not be presented with the same evidential strength as the leading V routes.

\begin{figure}[htbp]\centering
\includegraphics[width=\linewidth]{figures/replication.png}
\caption{All 96 selected routes: common-panel versus additional-family means (left), and noising versus restoration means (right). Dashed lines indicate equality. This is selected discovery consistency, not held-out confirmation or a false-discovery estimate.}
\end{figure}

\subsection{Controls support selectivity, not exclusivity}
The refinement contains 55 distinct control configurations: 17 other-site and 38 random-receiver controls, including the targeted L15H25 \texttt{is} controls. All controls have smaller mean absolute family effects than their linked target route. None of the random-receiver controls meets the retention rule. However, seven other-site controls do meet it. For example, the L17H1 other-mother to L18H24 V control has mean absolute effect $0.399$, compared with $0.548$ for the target mother route on the same core population.

This supports a relative preference for the selected sites/receivers under the tested interventions, not exclusive communication from the answer fact. Control channel norms were recorded, but random controls were not norm-matched; these comparisons are not a calibrated null distribution or significance test. Controls were targeted to high-ranked routes, not applied uniformly to all 769 configurations. Some control IDs are shared after deduplication, so counts should not be interpreted as 24 entirely disjoint control triplets.

\begin{figure}[htbp]\centering
\includegraphics[width=\linewidth]{figures/controls.png}
\caption{Paired target/control summaries on the same 20 families. All points lie below equality, but several other-site effects are substantial. Norm matching was not enforced.}
\end{figure}

\subsection{K/V and source sites are not uniformly additive}
Forty paired interaction contrasts can be computed from measured matching terms on the common panel; eight meet the operational rule. For L17H1 mother to L20H1, $I_{KV}-I_K-I_V=+0.277$ on average, opposing part of the negative V-only effect. For L17H1 \texttt{is} to L18H18, the signed interaction mean is $+0.096$, but the mean absolute family interaction is $0.180$.

For L17H1 to L18H18 V, the union-minus-sum-of-three-sites contrast has signed mean $+0.058$ but mean absolute family magnitude $0.293$. Its corresponding L18H19 values are $+0.091$ and $0.125$. Thus approximate average additivity would hide meaningful family-level interactions. The 136 unavailable contrast entries in the original comparison output reflect missing matched component measurements, including controls; they are not zeros, failed model runs, or 136 additional required GPU experiments.

\begin{figure}[htbp]\centering
\includegraphics[width=\linewidth]{figures/interactions.png}
\caption{The eight interaction contrasts meeting the frozen common-panel rule. Bars compare signed means with mean absolute family effects; they are not uncertainty bars. Neither a KV contrast nor a source-site union contrast is a measured multi-head interaction.}
\end{figure}

\subsection{Order and prefix sensitivity constrain the interpretation}
The leading L17H1 union-to-L18H18 V effect averages $4.657$ when the queried fact precedes the competing fact and $2.144$ when it follows. Its mean absolute within-family order difference is $2.528$ logits. For L18H19 V the corresponding means are $2.504$ and $1.645$, with absolute difference $1.046$. The effects remain positive, but their magnitudes are not order-invariant. This makes group-level order robustness an informative later test, not a redundant repeat of a solved question.

\begin{figure}[htbp]\centering
\includegraphics[width=\linewidth]{figures/order_sensitivity.png}
\caption{Each point pairs the two orders of one family for a fixed route. Consistent sign does not imply invariant magnitude. These analyses use existing S4.1 endpoints and required no new GPU computation.}
\end{figure}

Families 168 and 174 have shared answer prefixes. Every extended colon bypass has exactly zero effect on these families in both directions. This is the expected structural null: a frozen residual branch graph cannot move a colon-local residual change to the later metric position. It is not evidence against the source head. All-family and no-prefix summaries are retained, and other route types let later prefix computation remain live.

\subsection{Implementation scope versus the full protocol}
The implemented controls and extensions completed; no duplicate run is needed. Nevertheless, the broad protocol wording must be narrowed to the executed scope: per-configuration self-endpoint checks apply once per worker to head/MLP routes, not every pair or bypass; impossible routes are gate checks; wrong-site controls are omitted for Q and bypass; random receivers apply only to head routes; refinement controls are measured on 40 pairs, not automatically in reverse or on all 178. The additional-69 summaries and paired order comparisons in this report were computed locally from saved records. Named multi-step mediation, G1--G3, mean-ablation circuits and held-out validation remain unexecuted.
\FloatBarrier

\section{What S4.1 establishes about RI}
The map supports a meaningful intersection between RI candidates and task-relevant communication: L16H1/L16H21 receive L15H25 V influence and send colon Q influence; L18H19 is a major receiver and downstream source; L22H5 receives L17H1 V influence; L19H16 and L23H15 participate in opposing routes. L17H24 has a supported bypass and a local negative MLP route, while its L21H23 continuation is weaker. L26H23 and L26H31 are not settled by the current map.

The large non-RI source effects, particularly L17H1, L15H25 and the newly covered L13H18, remain important omissions of RI selection. These observations support a targeted conditional RI audit, not either blanket acceptance or blanket rejection of RI. Exact routes are only one part of the membership criterion: conditional removal/restoration tests and eventual frozen validation are still required. \file{data/ri31_route_ledger.csv} records the observed status without labelling untested candidates absent.

\section{S4.2: concrete prospective measurements}
\label{sec:s42}
\textbf{Status: proposed from S4.1 evidence; no S4.2 GPU results or implementation are claimed.} The executable package remains S4.1-only. The separate \file{data/s42_proposal.json} fixes the proposed rosters, masks, background rule and deduplicated node combinations for review and subsequent implementation.

\subsection{G1: eight route-supported groups, 49 unique node configurations}
\begin{table}[htbp]\centering\footnotesize
\begin{tabular}{lp{.23\linewidth}p{.49\linewidth}r}
\toprule
ID & Pattern & Members & Configs \\
\midrule
O1 & L13H10 / sentence & L16H1, L17H3, L18H18, L18H19 & 9 \\
O2 & L15H25 / child-last & L16H1, L16H21 & 3 \\
O3 & L17H1 / mother & L18H18, L18H19, L18H24, L19H16, L20H1, L22H5, L27H6 & 15 \\
O4 & L17H3 / colon & L18H18, L18H19 & 3 \\
I1 & L27H6 / Q & L18H18, L18H19, L19H16, L20H1, L23H15, L25H17 & 13 \\
I2 & L21H23 / KV & L17H24, L18H18, L18H19, L19H16, L20H1 & 11 \\
I3 & L30H18 / Q & L18H18, L18H19, L20H1, L23H15, L25H17 & 11 \\
I4 & L18H18 / Q & L16H1, L16H21, L17H3 & 7 \\
\bottomrule
\end{tabular}

\caption{Proposed G1 groups. Outgoing rows list receiver heads; incoming rows list source heads. All \emph{group members} are patched at the colon. Source fact masks in outgoing labels identify the evidence motivating the group, not a replacement of the receiver mask. Counts are unique singles, full group and group-minus-one within each row, before cross-row deduplication.}
\label{tab:groups}
\end{table}

Outgoing groups are taken in the protocol's source order. L14H26 contributes a singleton route and therefore does not create an artificial multi-member group. For L17H1, the mother mask is the primary common-source definition because all seven V routes were extended; the measured writer-site union supplies a complementary blocking mask, with the L27H6 union edge explicitly common-panel only. MLP recipients remain separate from the head groups.

Incoming groups follow the common-panel sum of mean absolute retained route effects, preserving receiver/site/channel and deduplicating identical node masks. The L21H23 KV and V rosters are identical, so they require one set of node-group interventions. The L18H18 and L18H19 Q incoming rosters are also identical; their shared source-group measurements are reused while receiver-specific claims remain distinct. I2 deliberately preserves L17H24 as a provisional member nominated by core KV evidence; its weak expanded V result prevents treating membership as settled. No post-hoc sign or magnitude threshold is introduced to remove it.

For each group $S$, measure singles, $S$, and $S\setminus\{h\}$ using simultaneous paired-donor head-output replacement. Compute
\[
 N(S)=I(S)-\sum_{h\in S}I(h),\qquad
 \delta_h=I(S)-I(S\setminus\{h\}).
\]
The second contrast measures the additional intervention effect of $h$ when its partners are already replaced. A head left unpatched remains live and can still be affected by upstream changes. A large contrast is evidence of conditional dependence, not automatically sufficient backup; task performance must support that stronger label.

The eight rows contain 72 within-group configurations in total, but only \textbf{49 distinct colon node masks}: 13 singles and 36 multi-head masks. This is 1,960 common-panel endpoint measurements before reuse; if the 13 saved Scope-F singles pass exact compatibility checks, only 1,440 new joint endpoints remain for this base noising panel. Seven singles are available in the bundled Stage-2 references; the other six must be sourced from the completed Stage-3 measurements, not treated as missing or replaced by S4.1 path effects. This count excludes blocking, new baselines, gates, optional sign splits, extensions and sensitivity runs. The overall G1 common-configuration cap remains 128, not 128 per group. Reuse requires identical inputs, metric, model, donor policy and mask; a path endpoint is not interchangeable with a whole-head endpoint.

For a retained nonadditivity or conditional contrast, first use the already available group-minus-one results. Only then add focused sign-based/pairwise splits within the remaining cap. Mixed groups O3, I1--I3 make cancellation especially plausible; signed averages alone must not decide whether to stop. Extend at most 16 informative group contrasts and their prerequisites to all 178 pairs, following the protocol's family-effect ranking and deterministic ties. Group membership is not inferred separately for every member from a single joint effect.

\subsection{Receiver blocking: specify the live-model communication test}
Start with four source-to-head-group patterns:
\begin{enumerate}
\item L13H10, whole query sentence, to O1.
\item L15H25, child-last and separately the whole query sentence, to O2.
\item L17H1, the measured mother/\texttt{is}/period union and separately the whole query sentence, to O3.
\item L17H3, colon, to O4. Do not substitute an unsupported early-writer mask.
\end{enumerate}
These are six source-mask patterns. For each, restore the source in a corrupted recipient, then repeat with the recipient group's colon outputs clamped to their original corrupted values. Include the group self-clamp without source restoration, and the reverse noising/clean-clamp conditions. The contrast is
\[
 B_{W,V}=M(x_r;W_c)-M(x_r;W_c,V_r).
\]
It tests whether the recipients' ability to respond is needed for source restoration. Report both $B$ and source restoration $J_W$; $B/J_W$, when stable, is an intervention ratio, not a unique mediated fraction. The three states in each direction give at most 36 endpoint configurations across the six patterns before deduplication and reuse. Including the 49 base G1 masks gives a conservative 85-configuration envelope, leaving room under 128 for justified splits. Gates and baseline capture forwards are additional computational costs.

L14H26 to L16H1 can receive a separate singleton blocking check without inventing a coalition. MLP14 for L13H10 and MLP19 for L18H18/L18H19 are separate mediator candidates, not silently pooled with heads. Joint L15H25+L17H1 restoration is \emph{not} an initial default: their tested direct consumers differ. The candidate convergence in Figure~\ref{fig:map} motivates later chain-specific blocking only after the basic group evidence is available.

\subsection{G2: retain the predefined RI audit and select K from actual measurements}
The original reduced roster remains 31 heads, with the exact list and four fixed layer-matched 23-head reference cohorts included in the proposal JSON. There is no new RI selection based on favorable S4.1 effects. Use these three RI groups:
\begin{itemize}
\item Positive strong: L16H1, L16H21, L17H24, L18H19, L22H5.
\item Negative strong: L19H16, L23H15, L26H23.
\item Residual: the remaining 23, including L26H31.
\end{itemize}

Choose $K$ from successive prefixes of (L21H18, L22H5, L21H6, L27H6), all patched at the colon. Measure these joint backgrounds; do not infer their effects by summing saved singles. Pick the smallest prefix leaving 30--80\% of the original gap; if none qualifies, use the closest fraction to 0.55 and flag floor effects. The unresolved incoming routes of L21H18/L21H6 are a research gap, not grounds to change this predeclared background list.

For each of the 31 candidates, compare $I(K\cup\{h\})-I(K)$ with its saved intact-model Scope-P effect. The candidate's mask is all original prompt positions, while $K$ is colon-only. If $h\in K$, use $K\setminus\{h\}$ as its background. Masks must be merged simultaneously; mixing Scope F and P singles would invalidate the comparison. There are 1,240 initial candidate endpoints plus distinct backgrounds. Test the three RI groups and four reference cohorts in intact and weakened backgrounds, giving at most 560 group/reference endpoints before reuse.

Extend candidates retained by their conditional effect or its paired change from the intact effect to 178 pairs. The eight strong RI candidates remain eligible for route follow-up even when a conditional average cancels. Repeat retained claims with mean replacement using consistently recomputed baselines. All discovery means must follow the frozen role/slot fallback and leave-one-family-out policy; no existing S4.1 donor measurement substitutes for that sensitivity experiment.

\subsection{G3: L26H31 remains an explicit backup question}
At the colon, measure L26H31 in four backgrounds: intact; L21H18 replaced; the triple L21H18, L22H5 and L21H6 replaced; and that triple plus L27H6 replaced. Record its attention/output before applying its own replacement, conditional answer effects, and reverse restoration for retained contrasts. G2's full-prompt candidate mask is not automatically reusable for this colon-only test. Add or reuse the L27H6-only attention-null control: as a downstream head, L27H6 cannot change L26H31's earlier attention, although it can change the latter's conditional answer effect. Attention resemblance alone does not establish backup.

\subsection{Execution order and efficient reuse}
Prepare and validate joint hooks/mean-mask machinery; verify reusable baselines and singles; run the 36 new G1 joint masks and the small K-selection panel; then measure the six blocking patterns, conditional RI singles/groups and nonduplicate G3 conditions. Freeze adaptive extensions only after complete-family verification. Cache compatibility must include the whole live/frozen background, not merely head IDs. Save both logits, candidate accuracy, intervention identities and necessary attention diagnostics once; derive interactions, order comparisons, bootstrap intervals and later compatible summaries on CPU.

L13H18 supplies a specific S4.3 local-mediation trigger that can be pursued alongside this plan when authorized. It does not require delaying all G1/G2/G3 work, and it does not justify a broad model-wide search immediately. S4.4/S4.5 should be designed jointly after these measurements, deduplicating identical model/input/mask configurations. New nonlinear interventions and adaptive circuit reductions will still require genuinely new model computation.
\FloatBarrier

\section{Conclusions and reserved later-stage results}
S4.1 supports a structured, task-specific communication map: strong fact-site V influence, colon Q continuation, parallel MLP/residual exits, and reproducible opposing contributions. The evidence is sufficiently concrete to nominate a small set of G1 groups and receiver-blocking patterns. It is not yet a faithful circuit, a universal semanticity result, or proof of individually confirmed RI membership.

The most material open questions are joint dependence versus compensation; whether the separately measured convergence motif actually composes; the source of order sensitivity; the unexplained L13H18 influence; and conditional contributions of the RI residual group and L26H31. These questions motivate the proposed S4.2 and bounded mediation work without repeating the completed S4.1 pipeline.

\subsection*{S4.2 results}
Not yet measured. The preceding section is the prospective plan to be reviewed before implementation and execution.
\subsection*{S4.3--S4.5 and frozen validation}
No results yet. These sections will be expanded in this same report after authorized experiments; the held-out population remains sealed.

\appendix
\section{Evidence files and reproduction}
Primary data: \file{results/stage4_all_v2/}; executed inputs/package:
\file{results/stage4_s41_v2/}; authoritative design:
\file{results/stage4_design_v2/manifest_proposal.json}; protocol:
\file{Stage4_Protocol.tex}. The Stage-3 interpretation is grounded in its completed report and independent reviews, not the superseded early Stage-4 proposal.

Local audit and supplementary analysis:
\file{results/stage4_review_20260924/verification.json},
\file{analyze_s41.py}, \file{build_report_assets.py},
\file{gate_errors.csv}, \file{all_events.csv.gz}, and \file{reproduced/}.
The package includes all figures and the following portable derived evidence:
\begin{itemize}
\item \file{data/route_ledger_96.csv}: complete extended-route ledger.
\item \file{data/extension_population_summary.csv}: 20/69/89-family summaries in both directions, including descriptive intervals.
\item \file{data/control_comparisons.csv}: paired control/target summaries.
\item \file{data/retained_not_extended.csv}: the 19 retained but unextended configurations.
\item \file{data/seed_map_coverage.csv}: tested and retained source/site/channel counts, including null panels.
\item \file{data/order_sensitivity.csv} and \file{data/prefix_summary.csv}: saved-data sensitivity analyses.
\item \file{data/ri31_route_ledger.csv}: route nominations without premature membership labels.
\item \file{data/s42_proposal.json}: proposed groups, exact RI/reference rosters and masks; not runnable GPU code.
\end{itemize}

For CPU reproduction from the project root, run the two supplied analysis scripts. The first re-verifies raw hashes/coverage and uses the source analyzer on saved chunks; the second builds the report tables and figures. No model weights or GPU are needed. The original protocol's bootstrap seed and retention thresholds are unchanged. Original GPU artifacts and analysis CSVs were not edited.
\end{document}
